\chapter{Various}

\section{Searching}
	\kactlimport{fstTrue.h}
	\kactlimport{lstTrue.h}
	\kactlimport{realTrue.h}
	\kactlimport{Ternary.h}

\section{Dynamic programming}

When doing DP on intervals: $a[i][j] = \min_{i < k < j}(a[i][k] + a[k][j]) + f(i, j)$, where the (minimal) optimal $k$ increases with both $i$ and $j$,
\begin{itemize}
\item one can solve intervals in increasing order of length, and search $k = p[i][j]$ for $a[i][j]$ only between $p[i][j-1]$ and $p[i+1][j]$.
\item This is known as Knuth DP. Sufficient criteria for this are if $f(b,c) \le f(a,d)$ and $f(a,c) + f(b,d) \le f(a,d) + f(b,c)$ for all $a \le b \le c \le d$.
\item Consider also: LineContainer (ch. Data structures), monotone queues, ternary search.
\end{itemize}

	\kactlimport{CircularLCS.h}
	\kactlimport{SMAWK.h}

\section{Debugging tricks}
	\begin{itemize}
		\item \texttt{signal(SIGSEGV, [](int) \{ \_Exit(0); \});} converts segfaults into Wrong Answers.
			Similarly one can catch SIGABRT (assertion failures) and SIGFPE (zero divisions).
			\texttt{\_GLIBCXX\_DEBUG} violations generate SIGABRT (or SIGSEGV on gcc 5.4.0 apparently).
		\item \texttt{feenableexcept(29);} kills the program on NaNs (\texttt 1), 0-divs (\texttt 4), infinities (\texttt 8) and denormals (\texttt{16}).
	\end{itemize}

\section{Optimization tricks}
	\subsection{Bit hacks}
		\begin{itemize}
			\item \texttt{x \& -x} is the least bit in \texttt{x}.
			\item \texttt{for (int x = m; x; ) \{ --x \&= m; ... \}} loops over all subset masks of \texttt{m} (except \texttt{m} itself).
			\item \texttt{c = x\&-x, r = x+c; (((r\^{}x) >> 2)/c) | r} is the next number after \texttt{x} with the same number of bits set.
			\item \texttt{ F0R(b,k) F0R(i,1<<K) if (i\&1<<b) D[i] += D[i\^{}(1<<b)]; } computes all sums of subsets.
		\end{itemize}

	\subsection{STL Utilities}
		\begin{itemize}
			\item \lstinline{ranges::rotate(p, ranges::min_element(p));} Lexicographically smallest cyclic shift.
			\item \lstinline{next_permutation(v.begin(), v.end());} Next permutation.
			\item \lstinline{prev_permutation(v.begin(), v.end());} Previous permutation.
			\item \lstinline{auto [mn, mx] = ranges::minmax(v);} Min and max in one pass.
			\item \lstinline{ranges::nth_element(v, v.begin()+k);} k-th smallest in $O(n)$.
			\item \lstinline{ranges::binary_search(v, x);} Check existence in sorted vector.
			\item \lstinline{ranges::lower_bound(v, x);} First element $>= x$.
			\item \lstinline{ranges::upper_bound(v, x);} First element $> x$.
			\item \lstinline{int cnt = ranges::upper_bound(v, R) - ranges::lower_bound(v, L);} Count of elements in range $[L,R]$.
			\item \lstinline{auto [l,r] = equal_range(v.begin(),v.end(),x);} Range of equal values.
			\item \lstinline{v.erase(unique(v.begin(),v.end()),v.end());} Remove duplicates (after sort).
			\item \lstinline{iota(v.begin(),v.end(),0);} Fill with increasing values.
			\item Prefix sum
			\begin{lstlisting}
				inclusive_scan(v.begin(), v.end(), pre.begin()+1,
				[](auto sum, auto x){ return sum + x; }, 0LL
				);
			\end{lstlisting}
		\end{itemize}
	\subsection{Pragmas}
	
		\begin{itemize}
			\item \lstinline{#pragma GCC optimize ("Ofast")} will make GCC auto-vectorize for loops and optimizes floating points better (assumes associativity and turns off denormals).
			\item \lstinline{#pragma GCC target ("avx,avx2")} can double performance of vectorized code, but causes crashes on old machines. Also consider older \lstinline{#pragma GCC target ("sse4")}.
			\item \lstinline{#pragma GCC optimize ("trapv")} kills the program on integer overflows (but is really slow).
		\end{itemize}

	% \kactlimport{BumpAllocator.h}
	% \kactlimport{SmallPtr.h}
	% \kactlimport{BumpAllocatorSTL.h}
	\kactlimport{FastIO.h}

\section{Other languages}
	% \kactlimport{Main.java}
	\kactlimport{Python3.py}
	\kactlimport{Decimal.py}
	% \kactlimport{Kotlin.kt}

\section{Miscellaneous}
	\kactlimport{GridDirection.h}